% Generated draft: requires mathematical and source review.
\documentclass[11pt]{article}
\usepackage[margin=25mm]{geometry}
\usepackage{fontspec}
\setmainfont{Latin Modern Roman}
\usepackage{amsmath,amssymb,mathtools,unicode-math}
\setmathfont{Latin Modern Math}
\usepackage{enumitem,needspace,xurl,hyperref}
\setlength{\parindent}{0pt}
\setlength{\parskip}{5pt}
\setlength{\emergencystretch}{3em}
\newcommand{\N}{\mathbb{N}}
\newcommand{\Z}{\mathbb{Z}}
\newcommand{\Q}{\mathbb{Q}}
\newcommand{\R}{\mathbb{R}}
\newcommand{\C}{\mathbb{C}}
\newcommand{\F}{\mathbb{F}}
\newcommand{\A}{\mathbb{A}}
\newcommand{\PP}{\mathbb{P}}
\newcommand{\E}{\mathbb{E}}
\newcommand{\eps}{\varepsilon}
\newcommand{\dd}{\,\mathrm{d}}
\begin{document}
\section*{UnsolvedMath ID: 800185}
\textbf{Source problem code:} catalogue entry
\subsection{Definitions and mathematical statement}
Let \(\kappa\) be a fixed real number with \(0<\kappa<1\). For a positive integer \(n\), let
\[
d(n)=\sum_{a\mid n,\ a\geq 1}1
\]
be its number of positive divisors. Interpret \(d_{\kappa}\) as the generalized divisor function: it is multiplicative, \(d_{\kappa}(1)=1\), and, for each prime \(p\) and integer \(r\geq1\),
\[
d_{\kappa}(p^r)=\frac{\kappa(\kappa+1)\cdots(\kappa+r-1)}{r!}.
\]
Equivalently,
\[
\sum_{n=1}^{\infty}\frac{d_{\kappa}(n)}{n^s}=\zeta(s)^{\kappa}\qquad(\operatorname{Re}s>1),
\]
where the power is defined by the Euler product. These are the conventional meanings of the divisor-function notation; the supplied record does not explicitly define them.

For each fixed \(\kappa\in(0,1)\), asymptotically evaluate
\[
S_{\kappa}(x)=\sum_{\substack{n\in\Z\\1\leq n\leq x}}d_{\kappa}(n)d(n+1)d(n+2)
\qquad\text{as }x\to+\infty.
\]
The requested evaluation should identify the leading asymptotic term, including its dependence on \(\kappa\). The supplied statement specifies neither a proposed main term nor a required error estimate or uniformity as \(\kappa\) varies.
\subsection{Short English statement}
For each fixed real \(0<\kappa<1\), determine the leading asymptotic behavior of the sum of \(d_{\kappa}(n)d(n+1)d(n+2)\) over positive integers \(n\leq x\). Here the divisor functions are interpreted in their conventional sense, which is not explicitly defined in the supplied record.
\subsection{Sources}
\begin{itemize}
\item Local UnsolvedMath registry, permanent ID \(800185\), exact source problem code \texttt{catalogue entry}; title: ``Moments of L-functions --- Moments of the Riemann zeta function --- Problem 37.'' The record attributes the problem to Hamieh, Kadiri, and Ng, with proposed year 2024.
\item Supplied source reference: Hamieh, Kadiri, and Ng, \emph{Moments of L-functions Problem List} (2024), Problem 37, \url{https://arxiv.org/abs/2405.17800}. The supplied provenance records access on 2026-10-16 and extraction by \texttt{arxiv\_latex}; these are supplied provenance statements, not an independently verified source visit.
\item The local registry labels the problem ``open.'' This is a supplied status label, not an independently verified assessment of the literature.
\end{itemize}
\end{document}
